\documentclass[10pt,a4paper]{amsart}
\usepackage[margin=19mm]{geometry}
\usepackage{amsmath,amssymb,mathtools}
\usepackage[scr=rsfs,cal=euler]{mathalfa}
\usepackage{xfrac}
\usepackage[unicode,hidelinks,bookmarks=false]{hyperref}
\newcommand{\ie}{i.e.\ }
\newcommand{\cf}{cf.\ }
\newcommand{\resp}{respectively\ }
\newcommand{\Subsection}{\subsection}
\providecommand{\nobreakdash}{\nobreak\hbox{-}}
\setlength{\emergencystretch}{2em}
\multlinegap0pt
\usepackage{bm}
\usepackage{bbm}
\usepackage{dsfont}
\usepackage{xspace}
\usepackage{cancel}
\usepackage{mathtools}
\usepackage{amsthm}
\makeatletter
\@ifundefined{Lemma}{\newtheorem{Lemma}{Lemma}}{}
\@ifundefined{Proposition}{\newtheorem{Proposition}{Proposition}}{}
\@ifundefined{Theorem}{\newtheorem{Theorem}{Theorem}}{}
\makeatother
\DeclareMathOperator{\divv}{div}
\title[Global well-posedness for the compressible Navier-Stokes equ]{Global well-posedness for the compressible Navier-Stokes equations with vacuum and smallness on coefficient-coupled scaling invariant quantity}
\author{Hao Xu, Xin Zhong}
\date{Source: arXiv:2606.16238v1 (CC BY 4.0)}
\begin{document}
\maketitle

\noindent\emph{Source and licence.} Global well-posedness for the compressible Navier-Stokes equations with vacuum and smallness on coefficient-coupled scaling invariant quantity. Version arXiv:2606.16238v1 (CC BY 4.0), distributed under the Creative Commons Attribution 4.0 International licence, as recorded in the arXiv licence metadata for this exact version and shown on the abstract page https://arxiv.org/abs/2606.16238v1. Full rights notice: licenses/arxiv-2606.16238-CC-BY-4.0.txt.\par\smallskip

\noindent\textit{Editorial scope.} Counted unit: the Proposition whose source label is \texttt{pro4.1} in the article (source0.tex lines 456--471, source label \texttt{pro4.1}), delivered together with the article's own complete original proof of it (source0.tex lines 472--659, one \texttt{proof} environment of 188 lines). No mathematics is written, repaired or re-derived: statement and proof are the article's own text line for line. Required context retained uncounted: 1.2 (lines 153--158); 2.1 (lines 220--224); 2.2 (lines 220--224); thm2.1 (lines 227--263); 3.4 (lines 392--394); 3.5 (lines 398--405); 3.7 (lines 398--405); lem3.4 (lines 409--418); lem3.5 (lines 423--439). Every definition, display and label of those lines is retained unchanged. Omitted from this package and not redistributed: 1--152, 159--219, 225--226, 264--391, 395--397, 406--408, 419--422, 440--455, 660--1139. Nothing inside a retained excerpt is omitted or altered: every retained body is delivered exactly as the article's lines are written, so no omission receipt of any kind accompanies this package. Citations: the retained excerpts carry no citation command at all, so no bibliography entry of the article is transcribed and no reference is redistributed. Reference adaptation: none was needed and none was made; every reference inside the delivered unit resolves inside it, so no unresolved reference or citation is produced. Printed slips kept literal: no printed word, symbol, hypothesis or number of the counted statement or of its proof was altered. Layout and class: the article's own class is \texttt{article} with packages including amsfonts, graphics, indentfirst, color, cite, latexsym, geometry, amsmath, amssymb, hyperref, epsfig, amscd, amsthm, mathrsfs; the wrapper is the lane's pinned \texttt{amsart} editorial wrapper at 10pt on A4, which restates the theorem-like environments the retained text needs and adds the article's own macro definitions that the retained text needs (\texttt{\textbackslash divv}), with extra wrapper packages (bm, bbm, dsfont, xspace, cancel, mathtools). The delivered numbering is the unit's own. Assets: no retained span contains a figure environment, a graphics-inclusion command, a PDF-page inclusion, a table or a TikZ picture; no image file is copied into this package, the package declares no asset dependency and no image file is redistributed with it. Licence: version arXiv:2606.16238v1 of this article is distributed under the Creative Commons Attribution 4.0 International licence, as recorded in the arXiv licence metadata for this exact version and shown on the abstract page https://arxiv.org/abs/2606.16238v1. A full rights notice is delivered at licenses/arxiv-2606.16238-CC-BY-4.0.txt. Characters: every retained excerpt is pure ASCII, so the delivered engine, XeLaTeX with its default Latin Modern text font, reports no missing character.
\begin{equation}\label{1.2}
\begin{cases}
\rho_t + \divv(\rho u) = 0,\\
(\rho u)_t +  \divv(\rho u \otimes u)+\nabla P= \mu \Delta u + (\mu + \lambda) \nabla \divv u.
\end{cases}
\end{equation}

\begin{gather}\label{2.1}
(\rho,\rho u)(x,0)=(\rho_0,m_0)(x), \ \  x\in \mathbb{R}^3,
\\ \label{2.2}
(\rho,u)\to(0,0) \ \ {\rm as} \ \  |x|\to\infty,  \ \  {\rm for}\ \ t\geq 0.
\end{gather}

\begin{Theorem}\label{thm2.1}
For any given $q\in(3, 6)$, assume that
\begin{equation}\label{2.3}
0 \leq\rho_0 \in L^{\gamma} \cap H^1 \cap W^{1,q},\ \
u_0 \in D^1.
\end{equation}
There exists a positive constant $\varepsilon$ independent of $\mu, \lambda, a, \gamma$, and the initial data such that if
\begin{equation}\label{2.4}
\mathbb{E}_0\triangleq \mu^{-5} \alpha^{6} \bar{\rho}^3 \mathbb{E}_1 (0) \mathbb{E}_2 (0) \leq \varepsilon,
\end{equation}
where
\begin{gather*}
\alpha \triangleq \max\{1 + a, \gamma\},\ \ \bar{\rho} \triangleq 1 + \|\rho_0\|_{L^{\infty }},\\
\mathbb{E}_1 (0) \triangleq \int \left(\frac{1}{2} \rho_0 |u_0|^2 + \frac{a}{\gamma - 1} \rho_0^{\gamma} \right) \mathrm{d} x,\\
\mathbb{E}_2 (0)
\triangleq \int \big[\mu |\nabla u_0|^2 + (\mu + \lambda)(\divv u_0)^2 + \left(2 \mu + \lambda \right)^{-1}P_0^2 \big] \mathrm{d} x,
\end{gather*}
then the problem \eqref{1.2}, \eqref{2.1}, and \eqref{2.2} have a unique global strong solution $(\rho, u)$ on $\mathbb{R}^3 \times (0,\infty)$ satisfying
\begin{equation}\label{2.5}
0\leq \rho(x,t)\leq \frac{3}{2}\bar\rho,\ \ {for \ all}\ \  (x,t)\in \mathbb{R}^3 \times(0, \infty)
\end{equation}
and
\begin{equation}\label{2.6}
\begin{cases}
\rho \in C([0, \infty);L^2) \cap L^\infty(0, \infty; L^{\gamma} \cap H^1 \cap W^{1,q}), \ \ \rho_t \in L^\infty(0, \infty;L^2\cap L^q),\\[1mm]
\sqrt \rho u \in C ([0, \infty);L^2),\ \ u \in L^\infty(0, \infty; D^1) \cap L^2 (0, \infty; D^1),\\[1mm]
\sqrt \rho \dot{u} \in L^2 (0, \infty; L^2), \ \ t \dot{u} \in L^2 (0, \infty; D^1),\ \ \sqrt{t \rho} \dot{u} \in L^\infty(0, \infty;L^2).
\end{cases}
\end{equation}

Moreover, there exists a positive constant $C$ depending on transport coefficients and initial norms such that, for any $t \geq 1$, $p \in [2,6]$, and $r \in (1,\infty )$,
\begin{equation}\label{2.7}
\|\nabla u\|_{L^p} \leq C t^{- 1 + \frac{1}{p}},\ \
\|P\|_{L^r} \leq C t^{- 1 + \frac{1}{r}},\ \
\|(\nabla F,\nabla \omega )\|_{L^2} \leq C t^{- 1}.
\end{equation}
\end{Theorem}

\begin{equation}\label{3.4}
\Delta F = \divv (\rho \dot u),\ \  \mu \Delta \omega = \nabla \times(\rho \dot u).
\end{equation}

\begin{gather}\label{3.5}
\|(\nabla F,\mu \nabla \omega )\|_{L^p} \leq C \|\rho \dot u\|_{L^p},
\\ \label{3.6}
\|(F, \mu \omega )\|_{L^p} \leq C \|\rho \dot u\|_{L^2}^{\frac{3p - 6}{2p}} \|(\mu \nabla u, (\mu + \lambda ) \divv u, P)\|_{L^2}^{\frac{6 - p}{2p}},
\\ \label{3.7}
\|\nabla u\|_{L^p}
\leq C \mu^{\frac{6 - 3p}{2p}} \|\nabla u\|_{L^2}^{\frac{6 - p}{2p}} \|\rho \dot{u}\|_{L^2}^{\frac{3 p -6}{2p}} + C \left( 2 \mu + \lambda \right)^{-1} \left( \|P\|_{L^p} + \|P\|_{L^2}^{\frac{6 - p}{2p}} \|\rho \dot{u}\|_{L^2}^{\frac{3 p -6}{2p}} \right).
\end{gather}

\begin{Lemma}\label{lem3.4}
 If $0\leq m,j\leq l$, then it has
$$
\|\nabla^{j}f\|_{L^p} \leq C\|\nabla^m f\|_{L^q}^{1-\theta} \|\nabla^l f\|_{L^r}^\theta,
$$
where $0\leq \theta\leq 1$ $($when $p=\infty$, it requires $0<\theta<1)$   and $j$ satisfies
$$
\frac{j}{3} - \frac{1}{p} = \left( \frac{m}{3} - \frac{1}{q} \right) (1-\theta) + \left(\frac{l}{3} - \frac{1}{r} \right) \theta.
$$
\end{Lemma}

\begin{Lemma}\label{lem3.5}
Suppose that the function $y$ satisfies
$$
y'(t) = g (y) + b'(t) \ \ on \ \ [0,T], \ \ y(0) = y^0,
$$
with $g \in C (\mathbb{R})$ and $y,b \in W^{1,1} (0,T)$. If $g(\infty ) = - \infty $ and
\begin{equation}\label{3.8}
b (t_2) - b(t_1) \leq N_0 + N_1 (t_2 - t_1),
\end{equation}
for all $0 \leq t_1 < t_2 \leq T$ with some $N_0 \geq 0$ and $N_1 \geq 0$, then $$
y(t) \leq \max \left\{y^0, \alpha_0 \right\} + N_0 < \infty\ \ on \ \ [0,T],
$$
where $\alpha_0$ is a constant such that
\begin{equation}\label{3.9}
g (\alpha) \leq - N_1 \ \ for \ \ \alpha \geq \alpha_0.
\end{equation}
\end{Lemma}

\begin{Proposition}\label{pro4.1}
Let the conditions in Theorem \ref{thm2.1} be satisfied.
There exists a positive constant $\varepsilon$ such that if
\begin{equation}\label{4.1}
\sup_{0 \leq t \leq T}\|\rho \|_{L^\infty } \leq 2 \bar \rho \ \ and \ \ \sup_{0 \leq t \leq T} \mathbb{E} (t) \leq \mathbb{E}_0^{\frac{1}{2}},
\end{equation}
then
\begin{gather}\label{4.3}
\sup_{0 \leq t \leq T} \mathbb{E}_1 (t) + \int^T_0 \left[ \mu \|\nabla u\|_{L^2}^2 + (\mu + \lambda ) \|\divv u\|_{L^2}^2 \right] \mathrm{d} t \leq \mathbb{E}_1 (0),
\\ \label{4.4}
\sup_{0 \leq t \leq T} \mathbb{E}_2 (t) + \int^T_0 \left[\|\sqrt{\rho} \dot{u}\|_{L^2}^2 + (2 \mu + \lambda )^{-2} \|P\|_{L^3}^3 \right] \mathrm{d} t \leq K \mathbb{E}_2 (0),\\
\label{4.2}
\sup_{0 \leq t \leq T}\|\rho \|_{L^\infty } \leq \frac{3}{2} \bar \rho, \ \ \sup_{0 \leq t \leq T}\mathbb{E} (t) \leq \frac{1}{2} \mathbb{E}_0^{\frac{1}{2}},
\end{gather}
provided that $\mathbb{E}_0 \leq \varepsilon.$
\end{Proposition}

\begin{proof}
{\bf Step I. Proof of \eqref{4.3}.}
The desired \eqref{4.3} follows from the standard elementary energy estimate.

{\bf Step II. Proof of \eqref{4.4}.} Taking the inner product of \eqref{1.2}$_2$ with $u_t$ in $L^2$, one has that
\begin{align}\label{4.5}
& \frac{1}{2} \frac{\mathrm{d}}{\mathrm{d} t} \int \left[ \mu |\nabla u|^2 + (\mu + \lambda )(\divv u)^2 \right] \mathrm{d} x + \int \rho |\dot{u}|^2 \mathrm{d} x\nonumber\\
& \quad = \int P \divv u_t \mathrm{d} x + \int \rho u \cdot \nabla u \cdot \dot{u} \mathrm{d} x\nonumber\\
& \quad = \frac{\mathrm{d}}{\mathrm{d} t} \int P \divv u \mathrm{d} x - \frac{1}{2 (2 \mu + \lambda )} \frac{\mathrm{d}}{\mathrm{d} t} \int P^2 \mathrm{d} x + \int \rho u \cdot \nabla u \cdot \dot{u} \mathrm{d} x \notag\\
& \qquad - \frac{1}{2\mu + \lambda } \int P u \cdot \nabla F \mathrm{d} x + \frac{\gamma - 1}{ \left(2 \mu + \lambda \right)} \int P \divv u F \mathrm{d} x\nonumber\\
& \quad \triangleq \frac{\mathrm{d}}{\mathrm{d} t} \int P \divv u \mathrm{d} x - \frac{1}{2 (2 \mu + \lambda )} \frac{\mathrm{d}}{\mathrm{d} t} \int P^2 \mathrm{d} x + \sum^3_{i = 1} I_i.
\end{align}
It deduces from (\ref{3.5}) with $p = 2$, (\ref{4.1}), and (\ref{3.7}) with $p = 3$ that
\begin{align*}
I_1
& \leq K \bar{\rho}^{\frac{1}{2}} \|\sqrt{\rho} \dot{u}\|_{L^2} \|\nabla u\|_{L^3} \|u\|_{L^6} \\
& \leq \frac{1}{8} \|\sqrt{\rho} \dot{u}\|_{L^2}^2 + K \mu^{-1} \bar{\rho} \|\nabla u\|_{L^2}^3 \|\rho \dot u\|_{L^2} + K \left(2 \mu + \lambda \right)^{-2} \bar{\rho} \|\nabla u\|_{L^2}^2 \left(\|P\|_{L^3}^2 + \|\rho \dot u\|_{L^2} \|P\|_{L^2} \right)\\
& \leq K \mu^{-5} \bar{\rho}^3 \|(\mu^{\frac{1}{2}} \nabla u, \left(2\mu + \lambda \right)^{-\frac{1}{2}} P)\|_{L^2}^2 \|\mu^{\frac{1}{2}}\nabla u\|_{L^2}^4 + \frac{1}{4} \|\sqrt{\rho} \dot{u}\|_{L^2}^2 + \frac{1}{16 (2 \mu + \lambda )^2} \|P\|_{L^3}^3,\\
I_2 + I_3
& \leq K \left(2 \mu + \lambda \right)^{-1} \gamma \|\nabla F\|_{L^2} \|P\|_{L^3} \|\nabla u\|_{L^2} \\
& \leq \frac{1}{4} \|\sqrt{\rho} \dot{u}\|_{L^2}^2 + K \left(2 \mu + \lambda \right)^{-2} \gamma^2 \bar{\rho} \|\nabla u\|_{L^2}^2 \|P\|_{L^3}^2\\
& \leq \frac{1}{4} \|\sqrt{\rho} \dot{u}\|_{L^2}^2 + \frac{1}{16 (2 \mu + \lambda )^2} \|P\|_{L^3}^3 + K \mu^{-5} \alpha^6 \bar{\rho}^3 \|\mu^{\frac{1}{2}} \nabla u\|_{L^2}^6.
\end{align*}
Inserting the above estimates into (\ref{4.5}), we get that
\begin{align}\label{4.6}
& \frac{\mathrm{d}}{\mathrm{d} t} \int \big[ \mu |\nabla u|^2 + (\mu + \lambda )(\divv u)^2 + \left( 2 \mu + \lambda \right)^{-1}|P|^2 \big] \mathrm{d} x + \int \rho |\dot{u}|^2 \mathrm{d} x\notag \\
& \quad \leq 2 \frac{\mathrm{d}}{\mathrm{d} t} \int P \divv u \mathrm{d} x + \frac{\|P\|_{L^3}^3}{4 (2 \mu + \lambda )^2} + K \mu^{-5} \alpha^6 \bar{\rho}^3 \mathbb{E}_2 (t) \|\sqrt{\mu}\nabla u\|_{L^2}^4.
\end{align}

Next, for positive constant $K_1 > 1 $, one has that
$$
2 \int P \divv u \mathrm{d} x
\leq \frac{2\mu + \lambda }{4} \|\divv u\|_{L^2}^2 + K_1 \frac{ \|P\|_{L^2}^2}{2 \mu + \lambda } \leq \frac{1}{4} \mathbb{E}_2 (t) + K_1 \frac{ \|P\|_{L^2}^2}{2 \mu + \lambda }.
$$
Moreover, we rewrite (\ref{1.2})$_1$ as
\begin{equation}\label{4.7}
P_t + \gamma P \divv u + u \cdot \nabla P = 0.
\end{equation}
Multiplying (\ref{4.7}) by $2 K_1 P \left(2 \mu + \lambda \right)^{-1}$ in $L^2$ and using the fact
$\divv u = \frac{F + P }{2 \mu + \lambda }$ gives
\begin{align*}
& \frac{K_1}{2 \mu + \lambda } \frac{\mathrm{d}}{\mathrm{d} t} \int P^2 \mathrm{d} x + \frac{K_1 \left( 2 \gamma - 1\right)}{\left(2 \mu + \lambda \right)^2} \int P^3 \mathrm{d} x\notag\\
& \quad = - \frac{ K_1 \left(2 \gamma - 1 \right)}{\left(2 \mu + \lambda \right)^2} \int P^2 F \mathrm{d} x \leq \frac{K_1 \left(2 \gamma - 1 \right)}{\left(2 \mu + \lambda \right)^2} \|P\|_{L^3}^2 \|F\|_{L^3}\notag\\
& \quad \leq \frac{K_1 \left(2 \gamma - 1 \right)}{4 (2 \mu + \lambda )^2} \|P\|_{L^3}^3 + K \frac{\left(2 \gamma - 1 \right)}{(2 \mu + \lambda )^2} \|((2 \mu + \lambda )\divv u, P)\|_{L^2}^{\frac{3}{2}} \|\rho \dot u\|_{L^2}^{\frac{3}{2}}\notag\\
& \quad \leq \frac{K_1 \left(2 \gamma - 1 \right)}{4 (2 \mu + \lambda )^2} \|P\|_{L^3}^3 + \frac{1}{4} \|\sqrt{\rho} \dot{u}\|_{L^2}^2 + K \frac{\alpha^{4} \bar{\rho}^3}{\left( 2 \mu + \lambda \right)^8} \|((2 \mu + \lambda )\divv u, P)\|_{L^2}^6,
\end{align*}
which together with \eqref{4.6} indicates that
\begin{align*}
& \frac{\mathrm{d}}{\mathrm{d} t} \left[\mu \|\nabla u\|_{L^2}^2 + \left(\mu + \lambda \right)\|\divv u\|_{L^2}^2 + \frac{K_1 + 1}{2 \mu + \lambda } \|P\|_{L^2}^2 \right] + \frac{3}{4} \|\sqrt{\rho} \dot{u}\|_{L^2}^2 + \frac{K_1 \left( 2 \gamma - 1\right)}{2 \left(2 \mu + \lambda \right)^2} \|P\|_{L^3}^3\\
& \quad \leq 2 \frac{\mathrm{d}}{\mathrm{d} t} \int P \divv u \mathrm{d} x + K \left( 2 \mu + \lambda \right)^{-3} \mu^{-5} \alpha^4 \bar{\rho}^3 \|P\|_{L^2}^6 + K \mu^{-5} \alpha^6 \bar{\rho}^3 \mathbb{E}_2 (t) \|(\sqrt{\mu}\nabla u, \sqrt{\mu + \lambda} \divv u)\|_{L^2}^4.
\end{align*}
Integrating the above inequality over $(0,T)$ with respect to $t$ leads to
\begin{align*}
& \sup_{0 \leq t \leq T} \mathbb{E}_2 (t) + \frac{3}{4}\int^T_0 \|\sqrt{\rho} \dot{u}\|_{L^2}^2 \mathrm{d} t + \frac{K_1 \left( 2 \gamma - 1\right)}{2 \left(2 \mu + \lambda \right)^2} \int^T_0 \|P\|_{L^3}^3 \mathrm{d} t\\
& \quad \leq K \mathbb{E}_2 (0) + \frac{1}{4} \sup_{0 \leq t \leq T} \mathbb{E}_2 (t) + K \left( 2 \mu + \lambda \right)^{-3} \mu^{-5} \alpha^4 \bar{\rho}^3 \int^T_0 \|P\|_{L^2}^6 \mathrm{d} t\\
& \qquad + K \mu^{-5} \alpha^6 \bar{\rho}^3 \int^T_0 \mathbb{E}_2 (t) \|(\sqrt{\mu}\nabla u, \sqrt{\mu + \lambda} \divv u)\|_{L^2}^4 \mathrm{d} t\\
& \quad \leq K \mathbb{E}_2 (0) + \frac{1}{2} \sup_{0 \leq t \leq T} \mathbb{E}_2 (t) + \frac{1}{4 \left(2 \mu + \lambda \right)^2} \int^T_0 \|P\|_{L^3}^3 \mathrm{d} t,
\end{align*}
where one has used
\begin{align*}
& K \left( 2 \mu + \lambda \right)^{-3} \mu^{-5} \alpha^4 \bar{\rho}^3 \int^T_0 \|P\|_{L^2}^6 \mathrm{d} t + K \mu^{-5} \alpha^6 \bar{\rho}^3 \int^T_0 \mathbb{E}_2 (t) \|(\sqrt{\mu}\nabla u, \sqrt{\mu + \lambda} \divv u)\|_{L^2}^4 \mathrm{d} t\\
& \quad \leq K \left( 2 \mu + \lambda \right)^{-3} \mu^{-5} \alpha^4 \bar{\rho}^3 \int^T_0 \|P\|_{L^2}^2 \|P\|_{L^1} \|P\|_{L^3}^3 \mathrm{d} t + K \mu^{-5} \alpha^6 \bar{\rho}^3 \sup_{0 \leq t \leq T}\big(\mathbb{E}_1 (t) \mathbb{E}_2(t)^2\big)\\
& \quad  \leq K \mu^{-5} \alpha^5 \bar{\rho}^3 \sup_{0 \leq t \leq T}\big(\mathbb{E}_1(t) \mathbb{E}_2(t)\big)\frac{1}{(2 \mu + \lambda )^2} \int^T_0 \|P\|_{L^3}^3 \mathrm{d} t + K \mathbb{E}_0^{\frac{1}{2}} \sup_{0 \leq t \leq T} \mathbb{E}_2 (t)\\
&  \quad \leq \frac{K_2 \mathbb{E}_0^{\frac{1}{2}}}{(2 \mu + \lambda )^2} \int^T_0 \|P\|_{L^3}^3 \mathrm{d} t + K_3 \mathbb{E}_0^{\frac{1}{2}} \sup_{0 \leq t \leq T} \mathbb{E}_2 (t) \leq \frac{1}{4} \sup_{0 \leq t \leq T} \mathbb{E}_2 (t) + \frac{1}{4 (2 \mu + \lambda )^2} \int^T_0 \|P\|_{L^3}^3 \mathrm{d} t,
\end{align*}
provided that
$$
\mathbb{E}_0 \leq \varepsilon_1 \triangleq \min \left\{1, (4 K_2)^{-2}, (4 K_3)^{-2}\right\}.
$$

{\bf Step III. Proof of (\ref{4.2}).}
For any given $(x,t) \in \mathbb{R}^3 \times [0,T]$ and $\delta>0$, let
$$
\rho^{\delta} (y,s) = \rho (y,s) + \delta \exp\left\{- \int^s_0 \divv u \left( X(\tau; x,t), \tau \right)\mathrm{d} \tau \right\} > 0.
$$
Here $X(s;x,t)$ is given by
$$
\begin{cases}
\frac{\mathrm{d}}{\mathrm{d} t} X (s;x,t) = u \left(X (s;x,t),s \right), \ \ 0 \leq s < t,\\
X(t;x,t) = x.
\end{cases}
$$
Thus, \eqref{1.2}$_1$ yields that
$$
\frac{\mathrm{d}}{\mathrm{d} s} \rho^{\delta} \left(X (s;x,t),s \right) + \rho^{\delta} \left(X (s;x,t),s \right) \divv u \left(X (s;x,t),s \right) = 0.
$$
This leads to
$$
Y'(s) = g(s) + b'(s),
$$
where
$$
Y (s) = \ln \rho^{\delta} \left(X (s;x,t),s \right), \ \ g(s) = - \frac{a \rho^{\gamma} \left(X (s;x,t),s \right)}{2 \mu + \lambda },
$$
$$
b(s) = - \frac{1}{2\mu + \lambda } \int^s_0 F \left(X (\tau;x,t),\tau \right) \mathrm{d} \tau.
$$

For the effective viscous flux $F$, one deduces from (\ref{3.4})$_1$ that
$$
F \left(X (t;x,\tau),\tau \right) = - \frac{\mathrm{d}}{\mathrm{d}
\tau} \left[(- \Delta)^{-1} \divv (\rho u) \right] + [u_i,\mathcal{R}_{ij} ] (\rho u_j)
$$
with $[u_i, \mathcal{R}_{ij}] = u_i \mathcal{R}_{ij} - \mathcal{R}_{ij} u_i$ and $\mathcal{R}_{ij} = \partial_i \left(- \Delta\right)^{-1} \partial_j$, which gives
\begin{align}
b(t_2) - b(t_1)
& = \frac{1}{2 \mu + \lambda } \left[(- \Delta)^{-1} \mathrm{div} (\rho u) (t_2) - (- \Delta)^{-1} \divv (\rho u) (t_1) \right] - \frac{1}{2 \mu + \lambda } \int^{t_2}_{t_1} [u_i,\mathcal{R}_{ij} ] (\rho u_j) \mathrm{d} \tau \notag\\
& \leq \frac{2}{2 \mu + \lambda } \sup_{0 \leq t \leq T} \| (- \Delta)^{-1} \divv (\rho u)\|_{L^{\infty }} - \frac{1}{2 \mu + \lambda } \int^{t_2}_{t_1} [u_i,\mathcal{R}_{ij} ] (\rho u_j) \mathrm{d} \tau \triangleq J_1 + J_2.\label{4.8}
\end{align}
Now we estimate $J_1$ and $J_2$ term by term
\begin{align}\label{4.9}
J_1
& \leq K \mu^{-1} \|(-\Delta)^{-1} \divv (\rho u)\|_{L^6}^{\frac{1}{2}} \|\nabla (-\Delta)^{-1} \divv (\rho u)\|_{L^6}^{\frac{1}{2}} \leq K \mu^{-1} \|\rho u\|_{L^2}^{\frac{1}{2}} \|\rho u\|_{L^6}^{\frac{1}{2}}\notag \\
& \leq K \mu^{- \frac{5}{4}} \bar{\rho}^{\frac{3}{4}} \|\sqrt \rho u\|_{L^2}^{\frac{1}{2}} \left( \mu \|\nabla u\|_{L^2}^2 \right)^{\frac{1}{4}} \leq K \mathbb{E}_0^{\frac{1}{8}}.
\end{align}
Before estimating $J_2$, we deduce from (\ref{3.7}) with $p = 6$ that
$$
\|\nabla u\|_{L^6}
\leq K \mu^{-1} \bar{\rho}^{\frac{1}{2}} \|\sqrt{\rho} \dot{u}\|_{L^2} + K \mu^{-1} a \|\rho^{\gamma}\|_{L^6},
$$
which combined with  the commutator estimates and Lemma \ref{lem3.4} leads to
$$
\begin{aligned}
\|[u_i,\mathcal{R}_{ij} ] (\rho u_j)\|_{L^\infty }
& \leq K \|[u_i,\mathcal{R}_{ij} ] (\rho u_j)\|_{L^3}^{\frac{1}{5}} \|\nabla [u_i,\mathcal{R}_{ij} ] (\rho u_j)\|_{L^4}^{\frac{4}{5}}\\
& \leq K \|u\|_{L^6}^{\frac{1}{5}} \|\rho u\|_{L^6}^{\frac{1}{5}} \|\nabla u\|_{L^6}^{\frac{4}{5}} \|\rho u\|_{L^{12}}^{\frac{4}{5}}\\
& \leq K \bar{\rho} \|u\|_{L^6}^{\frac{2}{5}} \|\nabla u\|_{L^6}^{\frac{4}{5}} \Big( \|u\|_{L^6}^{\frac{3}{4}} \|\nabla u\|_{L^6}^{\frac{1}{4}}\Big)^{\frac{4}{5}} \leq K \bar{\rho} \|\nabla u\|_{L^2} \|\nabla u\|_{L^6}\\
& \leq K \mu^{-1} \bar{\rho}^{\frac{3}{2}} \|\nabla u\|_{L^2} \|\sqrt{\rho} \dot{u}\|_{L^2} + K \mu^{-1} a \bar{\rho} \|\nabla u\|_{L^2} \|\rho^{\gamma}\|_{L^6}. \\
\end{aligned}
$$
This implies that
\begin{align*}
J_2
& \leq K \left(2 \mu + \lambda \right)^{-1} \mu^{-\frac{3}{2}} \bar{\rho} \left(\int^{t_2}_{t_1} \|\mu^{\frac{1}{2}}\nabla u\|_{L^2}^2 \mathrm{d} \tau \right)^{\frac{1}{2}} \left[\int^{t_2}_{t_1} \left(\bar{\rho} \|\sqrt{\rho} \dot{u}\|_{L^2}^2 + a^2 \|\rho^{\gamma}\|_{L^6}^2 \right)\mathrm{d} \tau \right]^{\frac{1}{2}}\notag\\
& \leq K \mathbb{E}_0^{\frac{1}{4}} + K \mu^{-\frac{3}{2}} \left(2 \mu + \lambda \right)^{-1} a \bar{\rho} \left( \int^{t_2}_{t_1} \|\mu^{\frac{1}{2}}\nabla u\|_{L^2}^2 \mathrm{d} \tau \right)^{\frac{1}{2}} \left(\int^{t_2}_{t_1} \|\rho^{\gamma}\|_{L^6}^2\mathrm{d} \tau \right)^{\frac{1}{2}}\notag\\
& \leq K \mathbb{E}_0^{\frac{1}{4}}  + K \mu^{-\frac{3}{2}} \left(2 \mu + \lambda \right)^{-1} a \bar{\rho}^{\frac{2 + \gamma}{2}} \sup_{0 \leq t \leq T} \mathbb{E}_2(t)^{\frac{1}{6}} \left[\int^{t_2}_{t_1} \left(\|\mu^{\frac{1}{2}}\nabla u\|_{L^2}^2 \right)^{\frac{2}{3}} \mathrm{d} \tau \right]^{\frac{1}{2}} \left(\int^{t_2}_{t_1} \|\rho^{\gamma}\|_{L^3}\mathrm{d} \tau \right)^{\frac{1}{2}}\notag\\
& \leq K \mathbb{E}_0^{\frac{1}{4}}  + K \mu^{-\frac{3}{2}} \left(2 \mu + \lambda \right)^{-\frac{2}{3}} a^{\frac{1}{2}} \bar{\rho}^{\frac{2 + \gamma}{2}}  \left( t_2 - t_1 \right)^{\frac{1}{2}} \sup_{0 \leq t \leq T} \mathbb{E}_2(t)^{\frac{1}{6}} \left(\int^{t_2}_{t_1} \|\mu^{\frac{1}{2}}\nabla u\|_{L^2}^2 \mathrm{d} \tau \right)^{\frac{1}{3}}\notag \\
& \quad \times\left[\frac{1}{(2 \mu + \lambda )^2} \int^{t_2}_{t_1} \|P\|_{L^3}^3 \mathrm{d} \tau \right]^{\frac{1}{6}} \leq K \mathbb{E}_0^{\frac{1}{4}} + K \mathbb{E}_0^{\frac{1}{3}} + \frac{a \bar{\rho}^{\gamma}}{2 \mu + \lambda } \left( t_2 - t_1 \right),
\end{align*}
which together with \eqref{4.8} and \eqref{4.9} shows that
$$
b(t_2) - b(t_1)
\leq K\mathbb{E}_0^{\frac{1}{8}} \Big(1 + \mathbb{E}_0^{\frac{7}{24}} \Big) + \frac{a \bar{\rho}^{\gamma}}{2 \mu + \lambda } \left( t_2 - t_1 \right) \leq K \mathbb{E}_0^{\frac{1}{8}} + \frac{a \bar{\rho}^{\gamma}}{2 \mu + \lambda } \left( t_2 - t_1 \right) \leq N_0 + N_1 \left( t_2 - t_1 \right).
$$
Here
$$
N_0 = K \mathbb{E}_0^{\frac{1}{8}}, \ \ N_1 = \frac{a \bar{\rho}^{\gamma}}{2 \mu + \lambda }.
$$

Recalling that
$$
g = g (Y) \triangleq - \left(e^{Y} - \delta \exp\left\{- \int^s_0 \divv u \left(X(\tau;x,t), \tau \right) \mathrm{d} \tau \right\} \right)^{\gamma} \left( 2 \mu + \lambda \right)^{-1},
$$
which implies that $g(\infty ) = - \infty $. Assume that
$$
\bar{Y}_{\delta} \triangleq \ln \left\{ \bar{\rho} + \delta \exp\left\{\int^T_0 \|\divv u(\cdot,\tau)\|_{L^{\infty }} \mathrm{d} \tau \right\} \right\}.
$$
Let $Y \geq \bar{Y}_{\delta}$, then
$$
\begin{aligned}
g(Y)
& = - \left(e^{Y} - \delta \exp\left\{- \int^s_0 \divv u \left(X(\tau;x,t), \tau \right) \mathrm{d} \tau \right\} \right)^{\gamma} \left( 2 \mu + \lambda \right)^{-1}\\
& \leq - \left(e^{Y} - \delta \exp\left\{\int^T_0 \|\mathrm{div} u (\cdot,\tau)\|_{L^{\infty }} \mathrm{d} \tau \right\} \right)^{\gamma} \left( 2 \mu + \lambda \right)^{-1} \leq - N_1.
\end{aligned}
$$
Consequently, one gets from Lemma \ref{lem3.5} that
$$
\rho^{\delta} (x, s) \leq \max\left\{ \bar{\rho} + \delta, \exp\{\bar{Y}_{\delta} \}\right\} \exp\{N_0\}.
$$
Letting $\delta \rightarrow 0$, one has
$$
\rho(x,s) \leq \bar{\rho} \exp\{N_0\} \leq \bar{\rho} \exp\left\{K_4 \mathbb{E}_0^{\frac{1}{8}}\right\} \leq \frac{3}{2} \bar{\rho},
$$
provided that
$$
\mathbb{E}_0 \leq \varepsilon_2 \triangleq \min \left\{\varepsilon_1,  \left(\frac{ \ln \frac{3}{2}}{K_4} \right)^8\right\}.
$$

Finally, it follows from (\ref{4.3}) and (\ref{4.4}) that
$$
\mathbb{E} (t) \leq K_5 \mathbb{E}_0 \leq \frac{1}{2} \mathbb{E}_0^{\frac{1}{2}},
$$
provided that
$$
\mathbb{E}_0 \leq \varepsilon \triangleq \min \left\{\varepsilon_2, (2 K_5)^{-2} \right\}.
$$
The proof of Proposition \ref{pro4.1} is complete.
\end{proof}

\end{document}
